\scanpage{4}
Soit \(C\) une \(C^*\)-Algèbre, on veut déterminer les \deleted{\unsure{endomorphismes}} endomorphismes linéaires continus \(T\) permutant aux translations gauches. Soit \(T\) un tel, on \unsure{notera|voit} que \(T''\) est un endomorphisme \unsure{fblt} continu \inserted{de \(C''\)} qui permute aux translations \(x\to xy\) \deleted{\((y\in C)\) \unsure{donc}} pour \(y\in C\), donc pour \(\forall y\in C''\) \inserted{donc} \(\bigl(T''(xy)=T''(x)y\quad\deleted{\unsure{il est même déjà}}\inserted{les deux}\text{ mbs sont \unsure{fblt} continus}\bigr)\).

Donc de la forme \(x\to ax\) \((a\in C'')\)\inserted{\unsure{bien déterminé car} \(a=\unsure{T''(1)}\)}, En fait, \(a\) appartient à l'ens.\inserted{\(\Gamma_g\)} des éléments \(a\) de \(C''\) tels que \(aC\subset C\). On voit donc que l'ensemble d'endomorphismes considéré est isomorphe, avec sa structure d'algèbre normée, à \(\Gamma_g\), \deleted{\unsure{T herm.}} \(a\) hermitien signifie \(x^*T(x)\) hermitien pour tout \(x\in C\); \(a\) positif signifie \(x^*T(x)\geq0\) pour \deleted{tout} \(x\in C\).

Soit \deleted{\unsure{g}} \(\Gamma\) l'ens des \(a\in C''\) tels que \(a\uline{C}\subset\uline{C},\ \uline{C}a\subset\uline{C}\). C'est une sous. \(C^*\)-Algèbre involutive. \unsure{T. él\textsuperscript{t}|Tout élément} hermitien de \(\Gamma_g\) est dans \(\Gamma\).

Supposons maintenant \(C\) plongé dans la \(C^*\)-Algèbre \(R\), et soit \(\Gamma_1\) l'ens des \(x\in R\) tels que \(xC\subset C,\ \unsure{Ca|Cx}\subset R\), c'est une sous \(C^*\)-algèbre de \(R\). Pour tout \(a\in\Gamma_1\), soit \(\varphi(a)\) l'élément de \(C''\) correspondant : l'opération \(T_a x=ax\) dans \(C\). Si \(a\) hermitien, \(\varphi(a)\) est hermitien, donc dans \(\Gamma\), donc pour \(\forall a\in\Gamma_1\), \(\varphi(a)\in\Gamma\); De plus, \(\varphi\) est évidemment une représentation d'algèbre, compatible avec l'involution (\unsure{puisque}

\mspage{5}
